\section{Introduction}

Legal case retrieval is a core task in legal information systems. In Task~1 of the Competition on Legal Information Extraction and Entailment (COLIEE), the system must automatically find the noticed cases for each query case: if a case was referenced by the query case, we treat it as a noticed case for that query. To test whether a model can truly infer the supporting precedents behind a judgment, citation information has been redacted from all case texts. The source corpus comes from the Federal Court of Canada, and the whole pipeline from query input to noticed-case prediction must therefore be fully automatic.

Our study starts from the observation reported in THUIR's COLIEE 2023 paper \cite{li2023thuir}. Table~\ref{tab:thuir-table3-excerpt}, excerpted from that paper \cite{li2023thuir}, shows that SAILER outperforms BM25 when both are constrained to 512 tokens, but unrestricted BM25 becomes stronger again than truncated SAILER. This suggests that the relative performance gap in legal retrieval may depend not only on model architecture, but also on how much of the judgment the model is actually allowed to read.

COLIEE 2026 is itself a long-document benchmark. Table~\ref{tab:length-stats} shows that the 7,708 cleaned full-document cases have an average length of 5,031.36 tokens and a median of 3,684.5 tokens; 56.11\% are shorter than 4096 tokens, 87.69\% are shorter than 8192 tokens, and 94.80\% are shorter than 12,288 tokens. This directly motivates our later chunk-based reranking design: if each document is split into at most three 4096-token chunks, then most cases can still be covered almost completely without increasing the encoder limit.

\begin{table}[t]
  \centering
  \small
  \begin{tabular}{lr}
    \toprule
    Statistic & Value \\
    \midrule
    Documents & 7,708 \\
    Mean tokens & 5,031.36 \\
    Median tokens & 3,684.5 \\
    Max tokens & 144,941 \\
    $<4096$ tokens & 4,325 (56.11\%) \\
    $<8192$ tokens & 6,759 (87.69\%) \\
    $<12,288$ tokens & 7,307 (94.80\%) \\
    \bottomrule
  \end{tabular}
  \caption{Length statistics of COLIEE 2026 cleaned full-document cases. With at most three 4096-token chunks per case, 12,288 tokens is the effective coverage limit for chunk-based reranking.}
  \label{tab:length-stats}
\end{table}

Table~\ref{tab:thuir-table3-excerpt} makes the truncation issue concrete \cite{li2023thuir}. When the maximum length is fixed at 512 tokens, SAILER reaches F1 of 0.1385 and BM25 only 0.1012. However, when BM25 is not length-limited, its F1 rises to 0.1541 and becomes stronger than truncated SAILER. Our own experiments go one step further: once the encoder context is extended to 4096 tokens, an embedding-based retrieval model (dense retrieval model) can surpass unlimited-length BM25. As shown in Table~\ref{tab:main-results-valid}, ModernBERT with domain-adaptive pretraining (DAPT) and supervised fine-tuning (SFT) reaches 0.1650 F1 on COLIEE 2026 validation, higher than BM25 at 0.1499.

\begin{table}[t]
  \centering
  \small
  \begin{tabular}{lcccc}
    \toprule
    Model & max\_length & P@5 & R@5 & F1 \\
    \midrule
    BM25($k_1=3,b=1$) & 512 & 0.0963 & 0.1067 & 0.1012 \\
    SAILER & 512 & 0.1315 & 0.1459 & 0.1385 \\
    BM25($k_1=3,b=1$) & -- & 0.1465 & 0.1625 & 0.1541 \\
    \bottomrule
  \end{tabular}
  \caption{Key rows excerpted from THUIR@COLIEE 2023 Table~3 \cite{li2023thuir}. ``--'' indicates unlimited length.}
  \label{tab:thuir-table3-excerpt}
\end{table}

Based on this observation, we adopt ModernBERT \cite{warner2024modernbert} as the long-context encoder and further apply domain-adaptive pretraining (DAPT), implemented as continued pretraining with masked language modeling (MLM), on U.S. case law. In practice, our single NVIDIA RTX 4080 SUPER 16GB GPU limits all formal experiments to 4096 tokens. In addition, under this hardware budget, the DAPT stage runs for about 60 hours and covers only 0.384 epoch, so the model still does not see the full 6.73 million-document U.S. case-law corpus.

This paper makes three main contributions. First, we re-evaluate the long-context hypothesis on COLIEE 2026 and show that query representation interacts strongly with available context length. Second, we show that stable legal embedding-based retrieval depends mainly on negative selection: static BM25 hard negatives nearly fail, dynamic negative sampling stabilizes training, and year-based scope filtering provides a smaller additional gain. Third, we integrate embedding-based retrieval, scope filtering, LightGBM learning to rank (LTR), and a post-processing stage whose key component is dynamic cut-off into a complete system that finally ranks 7th in the official COLIEE 2026 Task~1 results \cite{coliee2026task1results}.
